%!TEX root = ../../main.tex
% ============================================================
% 数学\线段图\六年级.tex
% 本文件由 main.tex 通过 \input 引入，请勿单独编译
% ============================================================


\section{解答：分数应用题——足球与篮球的比例}

\vspace{0.4cm}

\noindent\textbf{2. 学校买了一些足球和篮球，足球的个数占两种球总个数的 $\frac{3}{5}$。已知买了 18 个篮球，买了多少个足球？（先画出线段图，再解答）}

\vspace{0.8cm}

\begin{center}
\begin{tikzpicture}[>=Stealth, line width=1pt, scale=1.0] % 整体思路：把总数量看成一条被平均分成 5 份的线段，其中前 3 份表示足球，后 2 份表示篮球，再用大括号标出这两部分
  % 主线段 (10 units, 5 parts)
  \draw[thick] (0, 0) -- (10, 0); % 画总线段，长度 10 便于平均分成 5 份，每份长 2
  
  % 刻度 (6个刻度划分5个等份)
  \foreach \x in {0, 2, 4, 6, 8, 10} { % 在 0、2、4、6、8、10 处画刻度，把总段分成 5 等份
    \draw[thick] (\x, 0.2) -- (\x, -0.2); % 每个刻度都是一小段竖线，向上 0.2、向下 0.2
  } % 结束刻度循环
  
  % 足球大括号 (下方，前3份)
  \draw[decorate, decoration={brace, amplitude=8pt, mirror}] (0, -0.5) -- (6, -0.5) node[midway, below=12pt, font=\large] {足球}; % 在前 3 份下方画朝下的大括号，表示足球占 3/5
  
  % 篮球大括号 (上方，后2份)
  \draw[decorate, decoration={brace, amplitude=8pt}] (6, 0.5) -- (10, 0.5) node[midway, above=10pt, font=\large] {篮球 18个}; % 在后 2 份上方画大括号，标出篮球共 18 个

\end{tikzpicture}
\end{center}

\vspace{0.6cm}

{\small\color{gray}
\textbf{解析：}\\
1. **分析比例关系**：题目指出足球占总个数的 $\frac{3}{5}$，这意味着如果把总数看作 $5$ 份，足球占其中的 $3$ 份。\\
2. **确定篮球比例**：剩余的 $5 - 3 = 2$ 份就是篮球。\\
3. **建立数量联系**：已知篮球有 $18$ 个，对应这 $2$ 份。因此平均每份是 $18 \div 2 = 9$（个）。\\
4. **求足球数量**：足球占 $3$ 份，所以足球数量为 $3 \times 9 = 27$（个）。
}

\vspace{0.4cm}
\noindent\textbf{解答：}\\
\indent $1 - \frac{3}{5} = \frac{2}{5}$ \dots\dots\dots 篮球占总数的几分之几\\
\indent $18 \div 2 \times 3 = 27$（个）

\vspace{0.2cm}
\noindent\textbf{答：}买了 27 个足球。

\vspace{0.6cm}

{\small\color{blue!70!black}
\textbf{绘图基本原理与步骤：}\\
\textbf{1. 确定分率对应关系}：根据分母（5）确定总份数，通过总长度比例化映射得出图中线段总长（本图取 10cm）。\\
\textbf{2. 绘制等分线段}：利用 \texttt{TikZ} 循环绘制等距刻度线，视觉化展示“单位 1”被平均分成了 5 份。\\
\textbf{3. 分段标注}：使用大括号 (\texttt{brace}) 装饰器区分不同种类的球。上方标注已知具体数值的“篮球”（2份），下方标注未知的“足球”（3份），形成清晰的数形结合。
}

\newpage


